%=====================================================================
\chapter{Conclusion and Future Scope}
\label{ch:conclusion}
%=====================================================================

This chapter consolidates what the thesis has established across its four research objectives,
states the limits of the evidence explicitly, explains why two accuracy figures reported in
different chapters are not in conflict, sets out the research significance, and identifies the
work that remains.

\section{Consolidated Outcomes Against the Objectives}
\label{sec:concl-outcomes}

Table~\ref{tab:outcomes} summarises what each objective produced and the status of the evidence
behind it. The final column is stated deliberately, because the strength of a result and the
confidence that can be placed in it are separate properties and have been treated separately
throughout this work.

\begin{table}[!htb]
\centering
\caption[Objective-wise consolidated outcomes and evidence status]{Objective-wise consolidated
outcomes and the status of the supporting evidence}
\label{tab:outcomes}
\tabfont
\begin{tabular}{@{}L{0.9cm}L{3.9cm}L{4.6cm}L{4.4cm}@{}}
\toprule
\textbf{Obj.} & \textbf{Principal outcome} & \textbf{Key quantitative evidence} &
\textbf{Status and limits} \\
\midrule
RO1 & Four-stream fusion model with adaptive, then uncertainty-conditioned, source weighting
(Ch.~\ref{ch:fusion}--\ref{ch:dynfusion}) & MAE 0.018, $R^2$ 0.988, DA 84.7\% against 76.3\% for a
fixed-weight ensemble on identical inputs; both textual sources contribute 10.9\% of $R^2$
superadditively; dynamic weighting adds 22.7\% Sharpe and removes 18.3\% of drawdown &
Established. Index-level, one-day horizon; macro interpolation understates announcement-day
effects; dynamic result reported relative to a baseline \\
RO2 & Hybrid LLM sentiment and indicator framework in a reproducible environment
(Ch.~\ref{ch:llm}) & Sentiment is the highest-gain feature; correlation with price 0.04, with
momentum 0.74; accuracy 47.72\% against a 46.43\% baseline mean; highest prediction confidence;
buy-signal success 72\% and 65\% & Established but bounded. All models near the random classifier
at a single-name daily horizon \\
RO3 & Causal, multi-granularity transparency and a measured account of sentiment noise
(Ch.~\ref{ch:causal}) & Causal graph density 0.850 with an identifiable hub; correct external sign
in 16 of 16 series-by-period tests; market states separated at $p=1.8\times10^{-8}$ on a withheld
variable & Established for construct validity. Predictive content separately measured and not
usable; five instruments only \\
RO4 & Regime-conditional decision framework, ablated at matched risk and tested across markets
(Ch.~\ref{ch:allocation}) & Sharpe 0.878 against 0.628; drawdown $-18.98\%$ against $-36.08\%$;
wins in 18/18 specifications, 18/19 years, 4/4 stress episodes; holdout Sharpe 1.614 & Established
as consistent, not as statistically significant. Bootstrap interval $[-0.102, 0.574]$ includes
zero; signal layer returns a null \\
\bottomrule
\end{tabular}
\end{table}

Read as a ledger, the pattern is consistent. Adaptive fusion beats fixed-weight combination, but
index-level accuracy does not transfer to single-name daily direction. Language-model sentiment is
orthogonal and highest-importance, but does not overturn market efficiency at a daily horizon. The
sentiment construct is validated against four independent external series, but a validated
construct does not thereby forecast returns. Regime risk timing improves both the Sharpe ratio and
the drawdown at matched exposure, but the Sharpe gain is not statistically significant on eighteen
years of data. The defensive behaviour transfers to a second universe, but the framework does not
add risk-adjusted value in a concentrated single-factor one.

These bounding statements are not shortcomings appended to the results. They are the boundary of
the contribution, and stating them is what makes the positive claims credible.

\section{How the Contributions Compose}
\label{sec:concl-compose}

The four contributions are not four independent studies with a shared subject. Each was produced
in response to a limitation exposed by the one before it, and Figure~\ref{fig:integrated} shows the
composed system together with the two places where a finding redirected the design of the next
stage.

\begin{figure}[!htb]
\centering
\scalebox{0.80}{\begin{tikzpicture}[node distance=3mm]
\node[blkd,minimum width=100mm,minimum height=7mm] (i)
  {Heterogeneous inputs: market records, financial news, social media, macroeconomic series,
   volatility indices};

\node[blk,below=5mm of i,minimum width=100mm,minimum height=11mm] (c1)
  {\textbf{Ch.~4--5 --- integration.} Four-stream fusion with cross-modal attention; source
   weights conditioned on\\ measured, time-varying reliability \hfill
   \textit{establishes: what to fuse, and how much to trust each stream}};

\node[blkg,below=4mm of c1,minimum width=100mm,minimum height=11mm] (c2)
  {\textbf{Ch.~6 --- signal extraction.} Language-model sentiment fused with technical indicators
   in an optimised\\ architecture \hfill
   \textit{establishes: sentiment is the highest-importance feature, and its increment is
   bounded}};

\node[blka,below=4mm of c2,minimum width=100mm,minimum height=11mm] (c3)
  {\textbf{Ch.~7 --- transparency.} Directed causal structure replaces correlation; sentiment
   validated externally\\ and its noise measured \hfill
   \textit{establishes: which relationships are real, and how far sentiment can be trusted}};

\node[blkg,below=4mm of c3,minimum width=100mm,minimum height=11mm] (c4)
  {\textbf{Ch.~8 --- decision.} Regime-conditional allocation under cost, turnover and position
   constraints;\\ ablated at matched risk \hfill
   \textit{establishes: which layer produces the outcome, and where it transfers}};

\node[blkd,below=5mm of c4,minimum width=100mm,minimum height=7mm] (o)
  {An integrated framework whose every layer has been separately tested, with its applicability
   bounded by evidence};

\draw[ar] (i) -- (c1); \draw[ar] (c1) -- (c2); \draw[ar] (c2) -- (c3);
\draw[ar] (c3) -- (c4); \draw[ar] (c4) -- (o);
\draw[arb] (c2.east) -- ++(7mm,0) |- node[lbl,right=3mm,pos=0.25,align=left]
  {bounded increment\\$\Rightarrow$ measure noise} (c3.east);
\draw[arb] (c3.west) -- ++(-7mm,0) |- node[lbl,left=3mm,pos=0.25,align=right]
  {validated construct\\$\Rightarrow$ testable claim} (c4.west);
\end{tikzpicture}}
\caption[Integration of the four research contributions]{Integration of the four contributions.
Each layer is a response to a limitation exposed by the previous one, and the dashed paths show the
two places where a finding in one study redirected the design of the next.}
\label{fig:integrated}
\end{figure}

The logic of the progression is as follows. Chapter~\ref{ch:fusion} established that fusing four
sources improves prediction and that the value of each source is state-dependent, which raised the
question of how to weight a source whose reliability changes --- answered in
Chapter~\ref{ch:dynfusion} by deriving the weight from measured predictive uncertainty rather than
from learned average relevance. Chapter~\ref{ch:llm} isolated the language-model sentiment channel
and found it to be the most important single feature yet to deliver only a bounded increment in
raw accuracy, which raised the question of whether the limitation lay in the measurement or in the
horizon. Chapter~\ref{ch:causal} answered by separating the two: the sentiment construct was
validated externally and found sound, while its predictive content was measured separately and
found unusable, so the limitation is \emph{not} one of measurement. Chapter~\ref{ch:allocation}
followed directly. If the signal channel is bounded and the boundedness does not stem from poor
measurement, the framework must earn its value at the decision layer --- and the matched-exposure
ablation confirms exactly this: risk timing by market state works, and the signal layer does not.

The final system is therefore not the one originally envisaged, and it is better supported for
that reason. An architecture whose components have each been tested, one of which has been shown
not to work, is a more reliable engineering artefact than one whose aggregate performance conceals
which of its parts is doing the work.

\section{Reconciling the Divergent Accuracy Figures}
\label{sec:reconcile}

A reader comparing Chapters~\ref{ch:fusion} and~\ref{ch:llm} will notice an apparent contradiction:
the first reports 84.7 per cent directional accuracy, the second 47.72 per cent. The two numbers
are not in conflict, and understanding why is central to the significance of this research.

They answer different questions about different objects. Chapter~\ref{ch:fusion} predicts the
next-day level of a \emph{diversified index} from four fused sources. An index aggregates hundreds
of constituents, so idiosyncratic noise largely cancels and the remaining signal is dominated by
persistent market-wide state --- exactly the component that macroeconomic conditions, aggregate
sentiment and index-level technical structure describe well. Chapter~\ref{ch:llm} predicts the
next-day \emph{direction of an individual equity}, which carries its full idiosyncratic component,
has close to balanced classes, and offers no aggregation to average the noise away. Predictability
is therefore expected to be far lower, and this is exactly what the empirical asset pricing
literature reports~\cite{gu2020,fama1970}.

Three implications follow, and together they define the trajectory of the research programme.

\textbf{An accuracy figure is meaningless without its object and horizon.} Comparing 84.7 per cent
to 47.72 per cent as though they measured the same thing is a category error, and the fact that
both numbers appear in this thesis makes the point more sharply than any argument could.

\textbf{The level at which a system is applied should match where the signal lives.} Aggregate
market state is predictable; single-name daily direction is largely not. A system that attempts to
extract single-name daily alpha from the same information that describes aggregate state is
applying a tool at the wrong granularity.

\textbf{If the signal channel is weak where decisions are taken, value must be created at the
decision layer.} This is the most consequential implication and it is what Chapter~\ref{ch:allocation}
tests directly. The matched-exposure ablation confirms it: risk timing works where signal
selection does not.

The apparent contradiction between the two figures is therefore not an inconsistency to be
explained away. It is the finding that organises the thesis.

\section{Research Significance}
\label{sec:concl-signif}

\subsection{Methodological significance}

The work advances multi-source fusion from a question of \emph{what} to combine to a question of
\emph{how} to combine under changing reliability. Expressing fusion weights as a function of
measured, time-varying predictive uncertainty gives a source-weighting rule that self-corrects
within days rather than at the next retraining cycle, and does so in response to causes that need
not be represented in the training inputs. Distribution-aware normalisation, fitted on an expanding
window, prevents a heavy-tailed input from dominating the fused representation without leaking
future information into the scaler.

Separating the allocation decision into relative composition and total risk is a second
methodological contribution. It makes each layer independently falsifiable, and without it the
central ablation of this thesis could not have been constructed at all --- the matched-exposure
control requires that the level of risk be settable independently of its composition.

\subsection{Evaluation significance}

Three evaluation practices used here are uncommon in this literature and are, in the author's
judgement, the most transferable part of the work.

The first is \textbf{ablation at matched risk}, which separates the timing of risk from its level.
Comparing a de-risked strategy against a fully invested benchmark confounds the two and credits a
framework for merely holding less. The literature contains many risk-managed strategies whose
reported advantage has never been decomposed in this way.

The second is \textbf{disclosure of the entire specification sweep} rather than its best cell ---
the minimum required once multiple testing is taken seriously~\cite{harvey2015,bailey2014dsr,
valls2026}. In this thesis the reported specification sits near the middle of its sweep rather than
at the top, and reporting the range rather than the maximum costs roughly nine per cent of headline
Sharpe ratio. That cost is the price of a credible number.

The third is \textbf{separation of construct validity from predictive usefulness}. The sentiment
composite is shown to measure sentiment --- sixteen of sixteen correct signs against independent
external series, market states separated at $p = 1.8\times10^{-8}$ on a withheld variable --- and,
separately, shown not to forecast returns. Conflating those questions is a recurring error, and
keeping them apart changes what a null result means: it converts ``the measure is bad'' into ``the
predictability is absent'', which points research in a completely different direction.

\subsection{Engineering significance}

Two failure modes were diagnosed that generalise beyond this framework, and both are silent.

\textbf{Scale mismatch in signal-to-allocation conversion.} Passing an unscaled indicator into a
mean-variance objective lets the ratio of units between the linear and quadratic terms, rather than
the information in the signal, determine the solution, producing a degenerate nearly single-asset
portfolio. Information-ratio scaling~\cite{grinold2000} removes it. The diagnostic symptom is a
concentrated solution on almost every rebalance date.

\textbf{An inert risk layer.} An absolute volatility target calibrated for a concentrated equity
portfolio binds in one rebalance out of 224 on a diversified universe, so the layer is present,
consumes computation and does nothing. A relative budget is scale-free and remains active. The
diagnostic is a one-line counter of how often the constraint binds.

Mixture-model label anchoring has the same character: omitting it produces no error message, only a
silently scrambled regime series. The common property of all three is that the software runs
correctly and the results are plausible. This class of defect is under-reported precisely because
it leaves no trace, and documenting instances of it is a contribution in itself.

\subsection{Practical significance}

For a practitioner the results support a specific and limited prescription.

\begin{itemize}
\item \textbf{Fuse sources, but weight them by measured reliability} rather than by historical
average importance. The weighting rule, not the expert networks, produced the measured gain in
Chapter~\ref{ch:dynfusion}.
\item \textbf{Use language-model sentiment for what it demonstrably provides} --- orthogonal
information and better-calibrated confidence --- rather than for a directional edge it does not
provide at a daily horizon. Its buy-side signal quality is materially better than its sell-side
quality.
\item \textbf{Prefer causally selected features where the structure is recoverable}, because they
survive regime change better and, as Chapter~\ref{ch:causal} showed, retain relationships that a
correlational filter discards silently.
\item \textbf{Implement a regime risk budget in relative rather than absolute form}, evaluate it
against a volatility-matched benchmark, and expect it to help only where the universe contains
genuine dispersion in asset volatility.
\item \textbf{Size positions conservatively}, in the fractional form the growth-optimal criterion
recommends once estimation error is admitted~\cite{kelly1956}.
\item \textbf{Instrument every risk layer with a binding counter}, so that an inert mechanism
announces itself.
\end{itemize}

For the Indian market specifically, the measured weights matter. Macroeconomic indicators carry 25
per cent of average attention, reflecting the influence of policy rates, inflation and industrial
output on valuations. Sentiment sources jointly carry 40 per cent, with the social media share
rising from 18 to 28 per cent at market extremes --- corroborating independently published findings
on the asymmetric influence of social sentiment on Indian sectoral
returns~\cite{khan2025asymmetric}. Sectoral predictability varies systematically with information
availability and factor exposure, from $R^2$ above 0.99 in financials to below 0.96 in real estate.
These are measured properties of that information environment rather than assumptions imported from
elsewhere.

\section{Conclusion}
\label{sec:concl-main}

This research set out to integrate multi-source data with sentiment analysis and language models
in order to improve stock market decision making, and to determine which parts of such an
integration actually contribute to the outcome. Both aims are met, though not entirely in the
direction originally anticipated, and the divergence is itself the most useful result.

A four-stream fusion model integrating financial metrics, news sentiment, social media sentiment
and macroeconomic indicators was developed and evaluated on ten years of Indian market data across
ten sectoral indices. It attains a normalised mean absolute error of 0.018, a coefficient of
determination of 0.988 and a directional accuracy of 84.7 per cent against 76.3 per cent for a
fixed-weight ensemble over identical inputs, with stability confirmed by walk-forward evaluation
at $0.0016 \pm 0.0004$ normalised error across five sequential folds. Removing both textual sources
costs 10.9 per cent of explained variance, exceeding the sum of their individual removals of 5.3
and 3.6 per cent --- the quantitative signature of complementary rather than redundant information.
The fusion rule was then advanced from learned average relevance to measured instantaneous
reliability through a Bayesian formulation in which weights are obtained by a softmax over negative
predictive uncertainties, improving the Sharpe ratio by 22.7 per cent and reducing maximum drawdown
by 18.3 per cent against the strongest static-fusion baseline while the same experts fused
statically add nothing.

A hybrid framework coupling a transformer language model with fifteen conventional technical
indicators was implemented in a reproducible computational environment and evaluated under a strict
temporal protocol. Language-model sentiment emerged as the highest-gain feature in the fused space,
correlating 0.74 with momentum yet only 0.04 with price, establishing that it supplies orthogonal
information; it produced the highest prediction confidence of the models examined, the most stable
accuracy across sentiment regimes, and buy-signal success rates near 72 and 65 per cent. The gain
in raw accuracy is modest --- 47.72 per cent against a baseline mean of 46.43 per cent --- and all
models operate close to the random classifier at a single-name daily horizon. That limit is
reported rather than concealed, and it redirected the research toward the decision layer.

Transparency was established structurally rather than by post-hoc attribution. A directed causal
adjacency matrix recovered by conditional independence testing yields a graph of density 0.850 with
an identifiable causal hub whose outgoing influence systematically exceeds its incoming influence,
and demonstrates that strong directed influence can coexist with weak linear correlation --- the
precise justification for causal rather than correlational feature selection, since a correlational
selector would have discarded a genuine driver. Sentiment noise was quantified rather than assumed:
the composite carries the expected sign against four independently published series, is significant
against all four in first differences, holds the correct sign in all sixteen series-by-sub-period
combinations, and separates the identified market states at $p = 1.8 \times 10^{-8}$ on a variable
withheld from the fitting procedure, while its predictive content, measured separately, is not
usable. Keeping construct validity and forecasting power apart is a contribution in its own right.

The decision framework was evaluated over eighteen years on ten liquid multi-asset instruments
under a walk-forward protocol with transaction costs charged to the framework and not to its
benchmarks. It attains a Sharpe ratio of 0.878 against 0.628 for an equal-weighted benchmark and
reduces maximum drawdown from $-36.1$ to $-19.0$ per cent, outperforming in all four pre-selected
stress episodes, in 18 of 19 calendar years on drawdown and in all 18 cells of a fully disclosed
specification sweep; on a 665-session holdout that postdates every design decision it records a
higher Sharpe ratio and a smaller drawdown than all three benchmarks. Ablation at matched average
exposure locates the source precisely: timing risk by identified market state improves both the
Sharpe ratio and the drawdown with no significant change in mean return, whereas the signal fusion
layer contributes nothing detectable, for reasons the information coefficients make plain.
Cross-market transfer reproduces the defensive behaviour and the regime effect in a concentrated
mega-capitalisation universe while losing on risk-adjusted return, bounding applicability to
universes with genuine dispersion in asset volatility.

The overall conclusion is precise rather than expansive. Multi-source integration and
language-model sentiment demonstrably improve the \emph{description} of the market, and causal
structure demonstrably improves the \emph{transparency} of that description. At the horizon and
granularity at which decisions are actually taken, however, the measurable value of the system
arises at the allocation layer, through regime-conditional risk timing, rather than at the signal
layer. A framework whose every layer has been separately tested, and whose applicability has been
bounded by evidence, is a more useful engineering artefact than one whose aggregate performance
conceals which of its components is doing the work.

\section{Limitations of the Thesis}
\label{sec:concl-limits}

The limitations specific to each study are stated in the chapter that reports it. Six apply to the
thesis as a whole and are consolidated here.

\textbf{Horizon.} Every result is at a daily decision horizon. Nothing in this thesis speaks to
intraday microstructure or to execution.

\textbf{Direction and leverage.} Every portfolio is long-only and unlevered, and risk budgeting is
one-sided. The framework cannot express a short view, and its defensive behaviour is achieved
entirely by moving toward cash.

\textbf{Language coverage.} Sentiment extraction is English-only. A substantial share of Indian
retail participants consume financial news in regional languages, so the measured sentiment channel
observes part of the information reaching the market and its measured influence is correspondingly
a lower bound.

\textbf{Direct rather than adaptive transfer.} The cross-market evidence comes from applying the
framework without modification. No claim about transfer \emph{learning} is made.

\textbf{No attribution experiments.} Transparency rests on structure; Shapley additive explanations
and local surrogate models are reviewed and cited but not run, so no comparative claim between
structural and attributional transparency is supported.

\textbf{The allocation universe is not Indian.} The allocation study is conducted on instruments
that satisfy its data requirements, and it is presented as a test of the mechanism rather than as a
claim about Indian allocation.

\section{Future Scope}
\label{sec:concl-future}

Eight directions follow from the evidence, ordered by how directly they address a limitation
identified above.

\textbf{1. Textual sentiment in a universe with instrument-level news coverage.} The null reported
in Chapter~\ref{ch:allocation} is a null about a market-derived proxy, not about textual sentiment.
Applying the fusion hypothesis where a time-stamped, instrument-level news corpus exists would
convert a proxy null into a genuine test, and would connect the language-model pipeline of
Chapter~\ref{ch:llm} directly to the allocation layer of Chapter~\ref{ch:allocation}. The
information-coefficient evidence of Section~\ref{sec:alloc-ic} --- positive and marginally
significant on single names, indistinguishable from zero on broad instruments --- suggests this is
where the signal, if any, will be found~\cite{officioso2026,liu2024llmreview}.

\textbf{2. Bayesian neural experts for epistemic uncertainty.} The uncertainty estimate of
Chapter~\ref{ch:dynfusion} is a rolling error variance, which captures aleatoric variation.
Replacing the point-estimating experts with Bayesian neural models would additionally expose
epistemic uncertainty, producing a confidence measure that distinguishes a noisy environment from
an unfamiliar one~\cite{gal2016dropout,kendall2017,khuwaja2023}. The operational difference is
that the system could report low confidence \emph{before} accumulating errors rather than after,
which is exactly the behaviour required at the onset of an unprecedented event.

\textbf{3. Non-stationary debiasing.} The stationarity assumption in the multi-scale debiasing layer
fails at exactly the moments that matter, as the 77.67 per cent peak volatility during the 2024
Union Budget window shows. A formulation whose noise model is itself state-dependent --- estimated
separately within identified market states, or allowed to follow the volatility process --- is the
direct remedy, and it connects naturally to the regime identification already present in
Chapter~\ref{ch:allocation}.

\textbf{4. Scalable causal discovery.} Extending conditional independence testing from five
instruments to a full fifty-constituent index requires an algorithm whose cost does not grow
combinatorially. This is a prerequisite for causal transparency at production scale, and without it
the structural-transparency argument of Chapter~\ref{ch:causal} remains a demonstration rather than
a deployable
method~\cite{liu2026dynamic,spirtes2000}.

\textbf{5. Adaptation-based cross-market transfer.} The natural extension of the direct-transfer
result is genuine transfer learning: pre-training the regime model and the source-reliability
estimates on a data-rich market and fine-tuning them on a thinner one, so that an emerging-market
deployment inherits structure learned elsewhere instead of starting from its own limited history.
The transferability condition identified in Section~\ref{sec:alloc-transfer} --- dispersion in
asset volatility --- gives a testable criterion for when this should be expected to work.

\textbf{6. An Indian multi-asset allocation study.} The allocation framework requires a universe
with volatility dispersion and long clean history. Nifty sectoral indices, government securities,
gold and real-estate instruments together satisfy the first requirement, and the history available
is now long enough to include the 2008 and 2020 episodes. This would close the one structural gap
between the prediction chapters and the decision chapter of this thesis.

\textbf{7. Learned risk budgets and explicit state persistence.} The regime multipliers are fixed
at 1.00, 0.75 and 0.50. Learning them --- reinforcement learning being the natural candidate given
the sequential structure~\cite{kabbani2022drl} --- would adapt the budget to the universe, subject
to the overfitting controls this thesis argues are mandatory~\cite{harvey2015}. Replacing the
Gaussian mixture with a hidden Markov or jump formulation would represent regime persistence
directly rather than recovering it incidentally~\cite{luo2026regime,hamilton1989}.

\textbf{8. Regional-language and higher-frequency sentiment, and deployment engineering.} Sentiment
extraction is currently English-only and daily. A large share of Indian retail participants read
financial news in regional languages, and intraday resolution would expose incorporation dynamics
that daily aggregation averages away --- particularly the two distinct incorporation regimes
identified in Section~\ref{sec:sent-leadlag}. On deployment, inference cost is negligible; the
binding constraints are the latency of retrieving and scoring text at scale, feed rate limits,
coordinated manipulation of social signals at exactly the market extremes where they carry most
weight, transaction costs, and compliance with algorithmic trading regulation. These are treated
here as first-class engineering problems rather than as afterthoughts, because a framework that
cannot be operated is not a decision-support system.
